% c02 —— §3 末 + §4 格极化 K3 曲面（前半）（p004–p005 / 原书 p.130–131）
% 承 c01 末式 (3.2)。p004 首行「$\triangle$ の極多面体(polar polytope) を」为跨页新句，
% 本块从「$\Delta$ 的极多面体（polar polytope）」起句；p005 末止于 (4.14) 后的「とおく．」。

% ===== 原书 p.130 =====

$\Delta$ 的极多面体（polar polytope）定义为
\begin{equation}
  \check{\Delta}=\{\,v\in\check{M}\mid\text{对任意 }m\in\Delta\text{ 有 }\langle v,m\rangle\ge -1\,\} \tag{3.3}
\end{equation}
这里 $\check{M}=\Hom(M,\mathbb{Z})$ 是 $M$ 的对偶格。当 $\check{\Delta}$ 为格多面体时，称 $\Delta$ 为自反的（reflexive）。对扇 $\Sigma$，其 1 维锥的生成元的凸包称为扇多面体（fan polytope）。适当选取以 $\Delta$ 为扇多面体的扇，把与它相伴的环面簇记作 $\check{X}$。当 $\check{X}$ 具有光滑的反典范超曲面 $\check{Y}$ 时，称 $Y$ 与 $\check{Y}$ 为 Batyrev 镜像对（Batyrev mirror pair）[Bat94]。

现在设 $X$ 为 3 维，取反典范线性系 $|-K_{X}|$ 与 $|-K_{\check{X}}|$ 中非常一般的因子 $Y$ 与 $\check{Y}$。把 $H^{2}(Y;\mathbb{Z})$ 中由限制映射 $\iota^{*}\colon H^{2}(X;\mathbb{Z})\to H^{2}(Y;\mathbb{Z})$ 的像生成的原始子格记作 $M_{\Delta}$，并同样定义 $M_{\check{\Delta}}\subset H^{*}(\check{Y},\mathbb{Z})$。

% [Tutor Note: 猜想 3.1 题名中的 (8.6) 是 Dolgachev [Dol96] 原文的公式号，不是本篇编号，不要计入 tag 计数。]
\begin{conjecture}[Dolgachev {[Dol96, Conjecture (8.6)]}]\mbox{}\par
\begin{enumerate}[label=\arabic*.,leftmargin=2.2em]
  \item 存在某个格 $\check{M}_{\Delta}$，使得 $M_{\Delta}$ 在 $H^{2}(Y;\mathbb{Z})$ 中的正交补格 $M_{\Delta}^{\perp}$ 正交分解为 $U\perp\check{M}_{\Delta}$。
  \item 存在原始嵌入 $M_{\check{\Delta}}\subset\check{M}_{\Delta}$。
  \item 等号 $M_{\Delta}=\check{M}_{\Delta}$ 成立的充要条件是 $M_{\Delta}\cong\Pic Y$。
\end{enumerate}
\end{conjecture}

% [Tutor Note: 原书源码对同一记号混用 \triangle 与 \Delta 两条命令（(4.1) 作 \triangle，(4.3) 左端与 p006 的 (4.18) 作 \Delta），印出的字形几乎相同；译本统一作 \Delta，数学含义不变。]

\section{格极化 K3 曲面}

在 K3 曲面 $Y$ 的 2 次上同调群 $H^{2}(Y,\mathbb{Z})$ 上以杯积引入配对，所得的格称为 K3 格。K3 格的秩为 22、符号为 $(3,19)$，是偶幺模格，并且作为抽象格同构于 $L=E_{8}\perp E_{8}\perp U\perp U\perp U$。对 K3 曲面 $Y$，令
\begin{equation}
  \Delta(Y)=\{\delta\in\Pic(Y)\mid(\delta,\delta)=-2\} \tag{4.1}
\end{equation}
对满足 $[\mathcal{L}]=\delta\in\Delta(Y)$ 的线丛 $\mathcal{L}$，由 Riemann–Roch 定理有
% [Tutor Note: 原书 (4.2) 右端 $=1$ 系把 $(\delta,\delta)=-2$ 代入 $2+\frac{1}{2}(\delta,\delta)$ 所得，并非笔误。]
\begin{equation}
  h^{0}(\mathcal{L})+h^{0}(\mathcal{L}^{\vee})\ge 2+\frac{1}{2}(\delta,\delta)=1 \tag{4.2}
\end{equation}
于是 $\mathcal{L}$ 与 $\mathcal{L}^{\vee}$ 中必有一方具有非平凡截面，从而是有效的；
\begin{align}
  \Delta(Y)&=\Delta(Y)^{+}\coprod\Delta(Y)^{-}, \tag{4.3}\\
  \Delta(Y)^{+}&=\{\delta\in\Delta(Y)\mid\delta\text{ 是有效的}\}, \tag{4.4}\\
  \Delta(Y)^{-}&=-\Delta(Y)^{+}. \tag{4.5}
\end{align}

% ===== 原书 p.131 =====

格的正交群 $O(L)$ 中由 $\Delta(Y)$ 中的元素给出的镜射生成的子群 $W(L)$，真不连续地作用于
\begin{equation}
  V(Y)=\{x\in H^{1,1}(Y)\cap H^{2}(Y,\mathbb{R})\mid (x,x)>0\} \tag{4.6}
\end{equation}
的包含 Kähler 类的连通分支 $V^{+}$ 上。该作用的基本区域由
\begin{equation}
  C(Y)=\{x\in V(Y)^{+}\mid\text{对任意 }\delta\in\Delta(Y)^{+}\text{ 有 }(x,\delta)\ge 0\} \tag{4.7}
\end{equation}
给出，Kähler 锥为
\begin{equation}
  C(Y)^{+}=\{x\in V(Y)^{+}\mid\text{对任意 }\delta\in\Delta(Y)^{+}\text{ 有 }(x,\delta)>0\} \tag{4.8}
\end{equation}
（例如参见 [BHPVdV04, Corollary VIII.3.9]）。由 Lefschetz 定理有
\begin{equation}
  \Pic(Y)=H^{1,1}(Y)\cap H^{2}(Y;\mathbb{Z}) \tag{4.9}
\end{equation}
注意到这一点，定义
\begin{align}
  \Pic(Y)^{+}&=C(Y)\cap H^{2}(Y;\mathbb{Z}), \tag{4.10}\\
  \Pic(Y)^{++}&=C(Y)^{+}\cap H^{2}(Y;\mathbb{Z}) \tag{4.11}
\end{align}

设 $M$ 为符号 $(1,t)$ 的偶非退化格。
\begin{equation}
  V(M)=\{x\in M_{\mathbb{R}}\mid (x,x)>0\} \tag{4.12}
\end{equation}
选取式 (4.12) 的一个连通分支，称之为 $V(M)^{+}$。
\begin{equation}
  \Delta(M)=\{\delta\in M\mid(\delta,\delta)=-2\} \tag{4.13}
\end{equation}
选取式 (4.13) 的子集 $\Delta(M)^{+}$，使其满足
\begin{enumerate}[label=\arabic*.,leftmargin=2.2em]
  \item $\Delta(M)=\Delta(M)^{+}\coprod\Delta(M)^{-}$，其中 $\Delta(M)^{-}=\{-\delta\mid\delta\in\Delta(M)^{+}\}$；
  \item $\Delta(M)^{+}$ 对加法封闭（但对减法不封闭），
\end{enumerate}
并令
\begin{equation}
  C(M)^{+}=\{h\in V(M)^{+}\cap M\mid\text{对任意 }\delta\in\Delta(M)^{+}\text{ 有 }(h,\delta)>0\} \tag{4.14}
\end{equation}

% CHUNK c02 DONE pages p004–p005
